\subsection{Parameterwise axioms and a stable pointed envelope}
\label{sec:A2-review}
This subsection uses covariant coherent diagrams to draw arrows. If
$\mathbb D_G$ denotes the manuscript's presheaf convention, put
$\mathbb D(J)=\mathbb D_G(J^{\op})$. The upper-left span $\Psi$ in Q2
then becomes the lower-right, cartesian corner $\Phi=\Psi^{\op}$; right
Kan extension is along that cartesian corner. In the covariant convention,
right Kan extension from the upper-left span would instead put a terminal
object at the last vertex and would not describe general cartesian squares.
We make this variance change throughout the provisional-axiom comparison.
Write $1$ for the terminal object and $0$ for the initial object of
$\mathbb D(e)$; they are not assumed isomorphic.

\begin{claim}{A2.1}{Proved}{All the stated coherent axioms already pass to shifts}
\label{A2-shifts}
Let $I$ be small. Each of the following properties of $\mathbb D$ passes to
$\mathbb D_I$: the final Der~6; either one-sided implication in that axiom;
each of the three implications of Der~8, with its stated cartesian,
cocartesian, or bicartesian hypothesis; and either provisional Der~6 or
Der~$6'$ with coefficients arbitrary, constant, or locally constant. For
constant and locally constant coefficients the condition in a shift is in
the indexing-category direction, with coefficient objects in $\mathbb D(I)$.
For $I\ne\varnothing$ the converse holds in each case. In addition,
\[
 W_{\mathbb D_I}=W_{\mathbb D}\quad(I\ne\varnothing),
 \qquad W_{\mathbb D_\varnothing}=\operatorname{Mor}(\Cat).
 \tag{A2.1.1}
\]
Consequently there is no derivator for which one of these coherently stated
axioms holds but fails after a shift. The margin at 105/051 is a valid demand,
but it adds no axiom to these formulations.
\end{claim}
\begin{proof}
For a functor $u:A\to B$, parameterwise Kan extension gives the canonical
isomorphisms
\[
 i^*(\id_I\times u)_!\cong u_!i^*,\qquad
 i^*(\id_I\times u)_*\cong u_*i^* \qquad(i\in I).
 \tag{A2.1.2}
\]
These are Beck--Chevalley maps, not arbitrary choices of isomorphisms. The
shift is a derivator, and the assertion that a square in $\mathbb D(I\times
\square)$ is (co)cartesian is equivalent to that assertion for its restrictions
at every $i$. Indeed, write the comparison with the appropriate corner Kan
extension and use (A2.1.2) and Der~2. The results used here are
\cite[Theorem~1.31 and Proposition~2.6]{A2GrothPreprint}; their proofs use
base change for the projection Grothendieck (op)fibrations, so do not assume
that $\mathbb D$ comes from a model category.

Evaluation preserves constant coefficients. It also preserves local
constancy: an arrow in the indexing direction which is invertible in
$\mathbb D(I)$ is invertible after every $i^*$. Thus a provisional axiom for
$\mathbb D$, applied at each $i$, proves the same axiom for $\mathbb D_I$.
The fibration or opfibration indexing the provisional axiom is unchanged.
The same pointwise argument applies to each implication of Der~8. Conversely,
pull a coherent square, cube, or coefficient object back constantly along
$I\to e$ and then evaluate at one $i$; (A2.1.2) identifies the prescribed
Kan-extension comparisons. If $I$ is empty all values of the shift are terminal.

Here is a separate verification for the localizer, including its exact
coefficient quantifiers. Write $p:A\to e$, $q:B\to e$, and $f:A\to B$.
The constant-Hom definition and the left Kan adjunction identify
$f\in W_{\mathbb D}$ with invertibility, for every $M\in\mathbb D(e)$, of
\[
 \tau_f(M):p_!p^*M\longrightarrow q_!q^*M.
 \tag{A2.1.3}
\]
For the reverse implication use Yoneda with all $N\in\mathbb D(e)$:
$\Hom(q_!q^*M,N)\to\Hom(p_!p^*M,N)$ is exactly the constant-Hom
comparison. In the shift the coefficient is an arbitrary $M\in\mathbb D(I)$.
By (A2.1.2), evaluation of (A2.1.3) at $i$ is $\tau_f(i^*M)$.
Der~2 proves the forward inclusion in (A2.1.1); constant $I$-diagrams
and one evaluation prove its converse. This rechecks the precise adjunction
step in GALOIS A.1--A.2 rather than imposing an enhancement on an equivalence
of ordinary categories.
\end{proof}

\begin{claim}{A2.2}{Proved}{The two unrestricted provisional axioms are different halves}
\label{A2-provisional}
When the coefficient in the provisional Der~6 is arbitrary, its exact content
is that every cartesian square is cocartesian. The unrestricted dual
provisional Der~$6'$ says that every cocartesian square is cartesian. Together
they are the final Der~6. Either half can hold without the other, even for
strong pointed derivators defined on all small categories.
\end{claim}
\begin{proof}
After the variance change specified above, write the provisional Kan
extension as $j_*p_*$, where $j$ is the cartesian corner inclusion. Its image
consists of cartesian squares. Conversely choose $p=\id$ of the corner,
which is both a fibration and an opfibration. Since $j_*$ is fully faithful,
every cartesian square is obtained in this way with an arbitrary corner
coefficient. The proof for $j_!p_!$ is dual. This argument uses arbitrary
coefficients: replacing them by constant or locally constant coefficients
invalidates the last essential-surjectivity argument.

For an explicit separation take a nonzero field $k$ and the model category
of simplicial $k$-vector spaces, equivalently nonnegative chain complexes
under Dold--Kan, with homology equivalences. Its strong derivator has
\[
 \mathbb C(I)=\operatorname{Fun}(I,\Ch_{\ge0}(k))
                   [\text{objectwise quasi-isomorphisms}^{-1}].
 \tag{A2.2.1}
\]
Restriction is precomposition and Kan extensions are derived Kan extensions.
Disjoint unions give products, weak equivalences are detected at each object,
and the comma formulas compute these Kan extensions. The projective and
injective diagram model structures give their adjunctions. Coherent arrows
and maps between them are obtained by cofibrant replacements and mapping
cylinders, giving strongness. These are the model-category verifications of
Der~1--Der~5 in \cite[Theorem~1.34, Proposition~1.36 and Theorem~1.38]{A2GrothPreprint}.

The associated infinity category embeds as the connective objects in
$D(k)$, preserving homotopy colimits; a homotopy pullback internally is the
connective truncation of the stable pullback. Thus every homotopy pushout is
a homotopy pullback. But the square
\[
 \begin{tikzcd}
 0\ar[r]\ar[d]&0\ar[d]\\
 0\ar[r]&k[0]
 \end{tikzcd}
 \tag{A2.2.2}
\]
is cartesian and not cocartesian: its ambient stable pullback is $k[-1]$,
whose connective truncation is zero, whereas the pushout of its zero span
is zero. The dual derivator
$\mathbb C^{\op}(I)=\mathbb C(I^{\op})^{\op}$ satisfies Der~1--Der~5
by duality and gives the other separation. Neither example acquires the
missing implication after a nonempty shift, by A2.1.

The homotopy derivator of spaces satisfies neither half. The pushout of
$*\leftarrow S^0\to *$ is $S^1$, but the pullback of
$*\to S^1\leftarrow *$ is $\Omega S^1\simeq\Z$, not $S^0$. Conversely the
cartesian square with initial corner $S^0\times S^0$, two middle corners
$S^0$, and terminal corner $*$ is not cocartesian: its homotopy pushout is
the four-edge cycle $K_{2,2}\simeq S^1$. Thus the margin is correct to reject
the preceding assertions. Its further proposed necessity is separately
refuted in A2.8. These are homotopy pushouts and pullbacks, not ordinary
set-theoretic ones.
\end{proof}

\begin{claim}{A2.3}{Proved}{The final Der 6 does not force a zero object}
\label{A2-unpointed}
Fix a nonzero spectrum $E$. The strong derivator of the slice infinity
category $\Sp_{/E}$ satisfies the final Der~6 at every shape, but is not
pointed. Its constant-coefficient localizer is $W_{\Z}$.
\end{claim}
\begin{proof}
Choose a combinatorial stable model category of spectra and a fibrant
representative of $E$. Use the slice model category and define the values
on all small categories by
\[
 \mathbb D_E(I)=\operatorname{Fun}(I,\mathrm{Spectra}_{/E})
                     [\text{objectwise stable equivalences}^{-1}].
 \tag{A2.3.1}
\]
This specifies restriction and all diagram categories, not just its base
homotopy category. The same diagram-model argument as in A2.2 verifies
Der~1--Der~5. Equivalently it is the homotopy derivator of $\Sp_{/E}$;
\cite[Example~2.5]{A2GrothPontoShulman} gives the derivator verification
for a complete and cocomplete infinity category.

The forgetful functor from a slice creates colimits and limits indexed by
categories with weakly contractible nerve;
in particular it creates the homotopy pushouts and pullbacks under discussion.
The relevant comparisons in spectra are therefore the same comparisons
as in the slice. Spectra have identical cartesian and cocartesian squares,
and the forgetful functor reflects equivalences. This proves final Der~6.
The initial object is $0\to E$ and the terminal object is $\id_E$, which
are inequivalent because $E\ne0$.

For the assertion about $W$, the pointed-envelope argument of A2.4 applies.
Pointed objects in $\Sp_{/E}$ are diagrams $E\to X\to E$ with the specified
composite the identity. They are equivalent to spectra by the fiber of
$X\to E$, with inverse $V\mapsto(E\to E\oplus V\to E)$. This works for
coherent diagrams of every shape. Thus $W_{\mathbb D_E}=W_{\Sp}$.
The latter tests whether $\Sigma^\infty(BA)_+\to\Sigma^\infty(BB)_+$ is an
equivalence (test its action on the sphere). Its cofiber is connective, so
it vanishes if and only if its integral homology vanishes: successive
stable Hurewicz statements kill its homotopy groups starting in degree zero.
This is exactly $W_{\Z}$, including maps with empty source or target.
\end{proof}

\begin{claim}{A2.4}{Proved}{The final Der 6 forces a homology localizer without pointedness}
\label{A2-stable-envelope}
Let $\mathbb D$ be a strong derivator satisfying the final Der~6. Its pointed
envelope $\mathbb D_*$ is a strong stable derivator, and
\[
 W_{\mathbb D}=W_{\mathbb D_*}=R(L_{\mathbb D_*}),\qquad
 L_{\mathbb D_*}=\{V\in\Ho(\Sp):V\otimes Z=0
                              \text{ for every }Z\in\mathbb D_*(e)\}.
 \tag{A2.4.1}
\]
In particular $W_{\mathbb D}=R(T(W_{\mathbb D}))$. Neither accessibility of
the localizer, a model presentation of $\mathbb D$, nor an initial assumption
that $\mathbb D$ is pointed is required. When $\mathbb D$ is pointed, final
Der~6 is precisely stability in the usual sense.
\end{claim}
\begin{proof}
For every $I$ let $\mathbb D_*(I)$ be the full subcategory of
$\mathbb D(I\times[1])$ consisting of coherent arrows whose source is
terminal in $\mathbb D(I)$. This is the derived pointing construction, not
the ordinary undercategory $1/\mathbb D(I)$. Its existence as a derivator
is \cite[III.5, Proposition~5.1 and Lemma~5.2, pp.~41--42]{A2Heller}, or
\cite[Proposition~4.11]{A2CisinskiKan}. The latter proof constructs a
reflector by the homotopy pushout of an arrow $A\to B$ along $A\to1$;
its formulas give both Kan extensions and the comma comparisons. Thus the
construction applies to a bare derivator. More explicitly, for $u:I\to J$
and a pointed diagram $Z$ its Kan extensions have underlying objects
\[
 U(u_!^{\mathrm{pt}}Z)=1_J\amalg_{u_!1_I}u_!UZ,\qquad
 U(u_*^{\mathrm{pt}}Z)=u_*UZ,
 \tag{A2.4.1a}
\]
where the left expression is a homotopy pushout, $u_*1_I\cong1_J$ supplies
the right pointing, and the superscript $\mathrm{pt}$ denotes Kan
extensions in the pointed derivator. These formulas
follow from the reflector and the Kan adjunctions; the pointwise
Beck--Chevalley formulas verify Der~4. Strongness is inherited: lift
an arrow or a commutative square in the full subcategory using strongness
of $\mathbb D^{[1]}$. The lifted source diagram is terminal, because its
values are terminal and the unique map to the terminal diagram is an
isomorphism by Der~2.

Let $U:\mathbb D_*\to\mathbb D$ evaluate at the target. It creates right
Kan extensions. It also creates homotopy pushouts: the colimit of the
constant source-$1$ diagram over the three-object span is $1$, since that
span is contractible (indeed it has an initial object). The reflector therefore
does not change the source of that left Kan extension. These observations
show that a square in $\mathbb D_*$ is cartesian, or cocartesian, exactly
when its image under $U$ is. The object $\id_1$ is both initial and terminal
in $\mathbb D_*$, so this derivator is pointed. Final Der~6 consequently
makes it stable.

There is an adjunction of derivators
\[
 F:\mathbb D\rightleftarrows\mathbb D_*:U,
 \qquad F(M)=(1\longrightarrow1\amalg M),
 \tag{A2.4.2}
\]
with $F$ cocontinuous; this is \cite[Lemma~4.15]{A2CisinskiKan}.
It is conservative in the present situation. Indeed for any coherent
$f:M\to N$ the natural square
\[
 \begin{tikzcd}
 M\ar[r]\ar[d,"f"']&1\amalg M\ar[d,"1\amalg f"]\\
 N\ar[r]&1\amalg N
 \end{tikzcd}
 \tag{A2.4.3}
\]
is a homotopy pushout, by coproduct pasting. Final Der~6 makes it a pullback;
if $F(f)$ is invertible, so is $f$. Strongness supplies a coherent representative
of any arrow in $\mathbb D(e)$, and the same argument works in every shift.

We next prove the required retract statement, including its maps. If
$Z=(s:1\to UZ)$ is pointed, there is a canonical cofiber square in
$\mathbb D_*$ with top map $F(s):F1\to FUZ$, bottom row $0\to Z$, and
right map the counit $\epsilon_Z:FUZ\to Z$. Under $U$ its top row is
$1\amalg1\to1\amalg UZ$, its left map is the fold, and its right map
is $[s,\id]$. It is a homotopy pushout by cancellation of the identity
summand. Thus it gives the triangle
\[
 F1\xrightarrow{F(s)}FUZ\xrightarrow{\epsilon_Z}Z\longrightarrow\Sigma F1.
 \tag{A2.4.4}
\]
The map $r_Z=F(UZ\to1)$ satisfies $r_ZF(s)=\id$. Hence
$(r_Z,\epsilon_Z):FUZ\to F1\oplus Z$ is an isomorphism, by the split
triangle criterion. In particular $Z$ is a retract of $FUZ$. The maps in
(A2.4.4) and $r_Z$ are natural in $Z$; even an objectwise choice of the
resulting retract would suffice for the next argument.

Use the precise comparisons $\tau_f$ of (A2.1.3). Cocontinuity of $F$
identifies $F\tau_f(M)$ with $\tau_f(FM)$. If $f\in W_{\mathbb D_*}$,
conservativity of $F$ makes every $\tau_f(M)$ invertible. Conversely if
$f\in W_{\mathbb D}$, then $\tau_f(FM)$ is invertible for every $M$.
For any $Z\in\mathbb D_*(e)$, naturality and (A2.4.4) express
$\tau_f(Z)$ as an arrow retract of $\tau_f(FUZ)$, so it too is invertible.
This proves the first equality in (A2.4.1).

For clarity we specify the spectral theorem being used. The published
\cite[Theorem~A.11, p.~1317]{A2CisinskiTabuada} says that evaluation at the
sphere identifies cocontinuous morphisms from the derivator of spectra to
a triangulated derivator $\mathbb E$ with objects of $\mathbb E(e)$.
It gives, for each $Z$, an exact coproduct-preserving functor
$V\mapsto V\otimes Z$, coherently compatible with Kan extensions. Therefore
\[
 p_!p^*Z\cong\Sigma^\infty(BA)_+\otimes Z,\qquad
 \cofib(\tau_f(Z))\cong\Sigma^\infty C(Bf)\otimes Z.
 \tag{A2.4.5}
\]
The intersection of the kernels of these exact coproduct-preserving functors
is a localizing subcategory $L_{\mathbb D_*}$. Equation (A2.4.5) proves
$W_{\mathbb D_*}=R(L_{\mathbb D_*})$ for every small-category functor.
Finally the defining Galois correspondence gives $RTR=R$: explicitly
$T(R(L))\subseteq L$ gives $R(T(R(L)))\subseteq R(L)$, and the reverse
inclusion holds because each cofiber which defines $R(L)$ is among the
generators of $T(R(L))$. Thus the displayed localizer is $RT$-fixed.
\end{proof}

\begin{claim}{A2.5}{Proved}{The cube's third implication and its limits}
\label{A2-Der8}
In any derivator, Der~8's implication $b+c\Rightarrow a$ is automatic for
cartesian squares, and $a+c\Rightarrow b$ is automatic for cocartesian
squares. Under final Der~6 the remaining implication $a+b\Rightarrow c$
also holds, at every shift, including the nonpointed case. Without final
Der~6 it can fail even for a represented strong derivator. Requiring only
this third implication for bicartesian squares does not imply final Der~6
or reverse excision for $W_{\mathbb D}$, and need not make that localizer
$RT$-fixed.
\end{claim}
\begin{proof}
The first two implications are invariance of the specified homotopy limit
or homotopy colimit under objectwise isomorphisms. Under final Der~6 apply
$F$ of A2.4 to a coherent map of bicartesian squares. It preserves
pushouts and reflects isomorphisms. In the stable target the two squares
supply compatible Mayer--Vietoris triangles
\[
 FX_0\longrightarrow FX_a\oplus FX_b\longrightarrow FX_1
                    \longrightarrow\Sigma FX_0;
\]
this is \cite[Theorem~6.1]{A2GrothPontoShulman}. If the maps at $0$ and $1$
are isomorphisms, the middle map $Ff_a\oplus Ff_b$ is an isomorphism.
Each summand map is its arrow retract, so both are isomorphisms, and
conservativity of $F$ concludes. Thus the margin's assertion that it suffices
to ask for the third implication is correct when all squares in the test
are already bicartesian; the first two add no requirements there.

For failure of the third implication in general let $P$ be the pentagon
lattice with $0<a<b<1$ and $0<c<1$, with $c$ incomparable to $a,b$.
It is complete and cocomplete. Define $\mathbb D_P(I)=\operatorname{Fun}(I,P)$.
Ordinary pointwise Kan extensions verify Der~1--Der~4, and coherent and
incoherent arrows agree, verifying strongness. In a lattice a square is
cartesian exactly when its upper-left vertex is the meet of its two middle
vertices, and cocartesian exactly when its lower-right vertex is their join.
The map
\[
 (0,a,c,1)\longrightarrow(0,b,c,1)
 \tag{A2.5.1}
\]
of bicartesian squares has identity maps at $0,c,1$ and the noninvertible
map $a\to b$. It contradicts $a+b\Rightarrow c$.

For the insufficiency assertion use the two-element chain $P=\{0<1\}$.
Its bicartesian squares satisfy the third implication: if the coordinatewise
inequalities $x\le x'$, $y\le y'$ preserve both meet and join, they preserve
each coordinate (check the three possible meet/join pairs). But the square
$(0,0,0,1)$ is cartesian and not cocartesian. Moreover $W_{\mathbb D_P}$
consists precisely of functors whose source and target are either both empty
or both nonempty. Indeed the only empty constant-Hom set is
$\Hom(\operatorname{const}1,\operatorname{const}0)$ over a nonempty shape;
all the other constant-Hom sets are singletons. Take the two-point discrete
category covered by its two singleton sieves. The inclusion of either
singleton into the cover belongs to this $W$, whereas the opposite-edge
map $\varnothing\to e$ does not. This disproves reverse square excision.
It also gives a precise failure of $RT$-fixedness. The map from the discrete
two-point category to $e$ belongs to $W_{\mathbb D_P}$; its pointed
unreduced cofiber is $S^1$. Consequently $T(W_{\mathbb D_P})=\Ho(\Sp)$,
whereas $\varnothing\to e$ is a named map in
$R(T(W_{\mathbb D_P}))\setminus W_{\mathbb D_P}$.
The finite lattice calculations, including the complete list of tests,
are checked in \texttt{code/A2\_verify\_lattices.py}.
\end{proof}

\begin{claim}{A2.6}{Proved}{A slice classification when an infinity category is given}
\label{A2-enhanced-classification}
Let $\mathcal C$ be a complete and cocomplete infinity category. All its
pushout squares are pullback squares and conversely if and only if
$\mathcal C$ is equivalent to $\mathcal E_{/V}$ for a complete and cocomplete
stable infinity category $\mathcal E$ and an object $V\in\mathcal E$.
One may take $\mathcal E=\mathcal C_*$ and $V=F1$, where $F\dashv U$
is free pointing. The equivalence identifies their homotopy derivators,
not merely their base homotopy categories.
\end{claim}
\begin{proof}
The reverse direction is the pullback and colimit argument in A2.3.
For the forward direction the proof of A2.4, performed in infinity
categories, makes $\mathcal E=\mathcal C_*$ stable. Define
\[
 K:\mathcal C\longrightarrow\mathcal E_{/V},\qquad
 X\longmapsto(FX\xrightarrow{F(!)}F1),
\]
and, for $q:Z\to V$ in $\mathcal E$, set
\[
 H(q)=UZ\times_{UV}1,
 \qquad 1\xrightarrow{\eta_1}UF1=UV.
 \tag{A2.6.1}
\]
These are functors into the genuine infinity-categorical slice. The
naturality square for the adjunction unit
\[
 \begin{tikzcd}
 X\ar[r,"\eta_X"]\ar[d,"!"']&UFX\ar[d,"UF(!)"]\\
 1\ar[r,"\eta_1"']&UF1
 \end{tikzcd}
\]
is a pushout, since it replaces the $X$ summand in $1\amalg X$ by $1$.
It is consequently a pullback, proving $HK\simeq\id_{\mathcal C}$.

For the other composite apply $F$ to the pullback defining $H(q)$. That
square is also a pushout by hypothesis; $F$ preserves it, and stability
makes the result a pullback. The natural splitting from (A2.4.4) identifies
\[
 FUZ\simeq V\oplus Z,\qquad FUV\simeq V\oplus V,\qquad
 FUq=\id_V\oplus q,\qquad F\eta_1=(\id_V,\id_V).
\]
The last equality uses both $F(!)\,F\eta_1=\id_V$ and
$\epsilon_{F1}F\eta_1=\id_V$, the adjunction triangle identity.
The pullback of $\id_V\oplus q$ along the diagonal of $V$ is $Z$, with its
map to $V$ equal to $q$. Hence $KH\simeq\id_{\mathcal E_{/V}}$.
The constructions and comparisons are natural, proving an equivalence
of infinity categories and therefore at every diagram shape.
\end{proof}

\begin{claim}{A2.7}{Proved}{An exact characterization for bare strong derivators}
\label{A2-bare-characterization}
For a strong derivator $\mathbb D$, final Der~6 is equivalent to these three
conditions on its canonical free-point adjunction:
\[
 \mathbb D_*\text{ is stable},\qquad
 F:\mathbb D\to\mathbb D_*\text{ is conservative},\qquad
 F\text{ preserves homotopy pullbacks}.
 \tag{A2.7.1}
\]
Here stability includes pointedness, which $\mathbb D_*$ always has.
This characterizes the final axiom using a canonical stable envelope, without
requiring an infinity-category presentation or a derived slice over a
nonterminal object. Under these conditions $\mathbb D$ is pointed if and
only if $F1=0$ in $\mathbb D_*(e)$.
\end{claim}
\begin{proof}
Under final Der~6 the first two conditions were proved in A2.4. A pullback
square in $\mathbb D$ is also a pushout; its image under the cocontinuous
$F$ is a pushout in the stable $\mathbb D_*$, hence is a pullback. This proves
the third condition.

Conversely assume (A2.7.1). Since $F$ is a left adjoint of derivators, it
preserves pushouts. If $Q$ is a pushout square in $\mathbb D$, then $FQ$ is
a pushout, hence a pullback. The canonical comparison of $Q$ with its
cartesian completion becomes an isomorphism after $F$: this uses the
specified pullback-preservation mate, not merely the fact that some pullback
exists in the target. Conservativity of $F$ reflects that isomorphism.
Thus $Q$ is cartesian. Dually, if $Q$ is cartesian, $FQ$ is cartesian and
therefore cocartesian; applying $F$ to the cocartesian-completion comparison
and using its cocontinuity and conservativity proves that $Q$ is cocartesian.
These constructions are parameterwise, so the same characterization holds
in every shift.
Finally $F0=0$ because $F$ is cocontinuous. Thus $F1=0$ makes
$F(0\to1)$ invertible and conservativity gives $0\cong1$; the converse is
immediate.
\end{proof}

\begin{claim}{A2.8}{Refuted}{The replacement conjecture in the HOT margin}
\label{A2-HOT-margin}
The final suggestion in 105/052, ms.~16, is false: a cartesian square of
spaces whose final vertex is $*$ can also be cocartesian while neither
middle vertex is contractible. For nonempty connected CW middle vertices
$A,B$, the exact condition is that $A*B$ be contractible, equivalently that
$\widetilde C_*(A;\Z)\otimes_{\Z}^{\mathbf L}\widetilde C_*(B;\Z)$ be acyclic.
\end{claim}
\begin{proof}
The cartesian square over $*$ has initial vertex $A\times B$. The homotopy
pushout of its two projections is the unreduced join $A*B$. Since $A,B$
are connected, van Kampen makes this join simply connected. Its reduced
chains are quasi-isomorphic to
$\Sigma(\widetilde C_*(A;\Z)\otimes_{\Z}^{\mathbf L}\widetilde C_*(B;\Z))$.
This follows also by collapsing the contractible CW subcomplex
$A*\{b_0\}\cup\{a_0\}*B$ to identify the join with $\Sigma(A\wedge B)$ up
to homotopy. Hurewicz and Whitehead now prove the stated equivalence.

Take $A=B=K$, the finite Higman presentation complex of D.5
(Claim~\ref{D5}). That claim gives the complete four-generator,
four-relator presentation, proves $\pi_1K\ne1$, and proves integral
acyclicity without assuming asphericity. Thus neither middle vertex is
contractible, but $K*K$ is simply connected and acyclic, hence contractible.
The square $(K\times K,K,K,*)$ is both cartesian and cocartesian.

The finite certificate is particularly direct. With the cellular bases in
D.5, $d_2:C_2K\to C_1K$ is a signed permutation matrix with determinant
$-1$. Its inverse defines a contraction $h:C_1K\to C_2K$ of reduced chains.
The product complex has chain ranks $(1,8,24,32,16)$. If $P$ is projection
to the degree-zero vertex, the integral homotopy
$h\otimes\id+P\otimes h$ satisfies $dH+Hd=\id-P\otimes P$.
The reduced join model has ranks $16,32,16$ in degrees $3,4,5$, and
the suspension of $h\otimes\id$ contracts it. The program
\texttt{code/A2\_verify\_lattices.py} constructs these maps and verifies the
identities on every basis element using integers. This is an integral
contraction certificate, not merely a rank calculation over a field.
\end{proof}

\begin{claim}{A2.9}{Proved}{Either unrestricted half already forces an RT-fixed localizer}
\label{A2-half-RT}
Let $\mathbb D$ be a strong derivator. If every homotopy pushout square is
cartesian, or if every homotopy pullback square is cocartesian, then
$W_{\mathbb D}$ is $RT$-fixed. Thus, for this conclusion about the
constant-coefficient localizer, either unrestricted provisional axiom of
A2.2 suffices. The opposite half, final Der~6, and Der~8 are not additional
requirements. The proof here uses unrestricted coefficients. Restricted-coefficient
versions are treated separately in A2.10--A2.14.
\end{claim}
\begin{proof}
First suppose that every pushout is a pullback. Put
$\mathbb E=\mathbb D_*$ and let $F:\mathbb D\to\mathbb E$ be free pointing.
The construction and Kan-extension formulas in A2.4 do not require
stability. Its forgetful functor creates pullbacks and pushouts, so
$\mathbb E$ has the same one-sided property. In a pointed derivator the
suspension square is then cartesian; hence the unit
$\id\to\Omega\Sigma$ is invertible. In particular $\Sigma$ is fully faithful.
The pushout square (A2.4.3) is also cartesian under the present one-sided
hypothesis, so $F$ is conservative. It is cocontinuous as before.
We do not assume that $\mathbb E$ is stable.

We give the representability argument in detail, so that no existence
theorem for a stabilization of $\mathbb E$ is needed. A pointed derivator
has the canonical cocontinuous action of pointed spaces
\[
 \Ho(\cS_*)\times\mathbb E(e)\longrightarrow\mathbb E(e),
 \qquad (X,P)\longmapsto X\otimes P.
\]
This follows from \cite[Corollary~4.19 and Theorem~5.22]{A2CisinskiKan};
it applies to a bare derivator. For $P,N\in\mathbb E(e)$ and $n\in\Z$
define
\[
 h^n_{P,N}(X)=
 \Hom_{\mathbb E(e)}
   \bigl(\Sigma^r(X\otimes P),\Sigma^{r+n}N\bigr),
 \qquad r\geq\max\{2,2-n\}.
 \tag{A2.9.1}
\]
All transition maps obtained by increasing $r$ are bijections, since
$\Sigma$ is fully faithful. Thus (A2.9.1) is canonically independent of the
choice of $r$. The double-suspension pinch, transported through the
pointed-space action, makes these sets abelian groups. Increasing $r$
preserves their group laws. Homotopy invariance is built into the action,
and reindexing $r$ gives natural suspension isomorphisms
$h^n(X)\cong h^{n+1}(\Sigma X)$.

The groups satisfy the wedge and cofiber exactness axioms on all pointed
CW complexes. For wedges this follows from cocontinuity of
$X\mapsto X\otimes P$ and of suspension, and the fact that maps out of a
coproduct form a product. The same single $r$ in (A2.9.1) works for an
arbitrary small wedge, so no interchange of a filtered colimit with a
product is being assumed.
For exactness, a pointed cofibration $A\to X$ yields, after tensoring and
suspending, a coherent cofiber sequence $L\to M\to C$ in $\mathbb E$.
If $g:M\to N'$ has zero composite with $L\to M$, strongness lifts the
commutative square from that coherent arrow to the coherent arrow
$0\to N'$. Applying the cofiber functor supplies $C\to N'$ whose
composite with $M\to C$ is $g$. Conversely every such composite kills
$L$. This proves
\[
 h^n_{P,N}(X/A)\longrightarrow h^n_{P,N}(X)
                    \longrightarrow h^n_{P,N}(A)
 \quad\text{exact}.
 \tag{A2.9.2}
\]
Together with the suspension isomorphisms this is a reduced generalized
cohomology theory.

Brown representability consequently supplies a spectrum $E_{P,N}$ with
\[
 h^n_{P,N}(X)\cong
 [\Sigma^\infty X,\Sigma^nE_{P,N}].
 \tag{A2.9.3}
\]
Here the representability is the all-CW, wedge-axiom theorem, with no
countability restriction: \cite[Theorem~II, condition~(4.6), and the
1963 correction]{A2Brown}, equivalently
\cite[Theorem~4E.1, pp.~448--452]{A2Hatcher}. The correction explains why
the full wedge axiom supplies the sequential-union condition. All Hom
groups and wedges are understood in the working universe, as for
the derivator's small Kan extensions.

It suffices to use $P=FM$, with $M\in\mathbb D(e)$, while $N$ ranges over
all of $\mathbb E(e)$. Set
\[
 L=\bigcap_{\substack{M\in\mathbb D(e),\,N\in\mathbb E(e)\\ n\in\Z}}
       \{V\in\Ho(\Sp):[V,\Sigma^nE_{FM,N}]=0\},
 \tag{A2.9.4}
\]
This is a localizing subcategory:
mapping into a fixed spectrum is cohomological, turns coproducts into
products, and all integer shifts occur in the test.
For $f:A\to B$, cocontinuity identifies
$F\tau_f(M)$ with $\tau_f(FM)$, and conservativity gives
\[
 \begin{aligned}
 f\in W_{\mathbb D}
 &\ \Longleftrightarrow\
 \tau_f(FM)\text{ is invertible for every }M\\
 &\ \Longleftrightarrow\
 h^n_{FM,N}((Bf)_+)\text{ is invertible for all }M,N,n.
 \end{aligned}
 \tag{A2.9.5}
\]
For the last equivalence, degree zero is ordinary Yoneda, since
$h^0_{P,N}(X)\cong\Hom(X\otimes P,N)$; the other degrees test the same
comparison against $\Sigma^nN$ for $n\geq0$, or against
$\Omega^{-n}N$ for $n<0$. The identification
$(BA)_+\otimes P\cong p_!p^*P$ is the canonical constant-hocolimit
formula. The long exact sequence of the cofiber of $(Bf)_+$ now identifies
(A2.9.5) with $\Sigma^\infty C(Bf)\in L$. Hence
$W_{\mathbb D}=R(L)$, and $RTR=R$ proves the conclusion.

Finally, if every pullback in $\mathbb D$ is a pushout, apply the
established case to $\mathbb D^{\op}$. Directly from constant Hom sets,
$f\in W_{\mathbb D^{\op}}$ if and only if
$f^{\op}\in W_{\mathbb D}$. Reversing an indexing category leaves its
classifying-space homotopy type unchanged, naturally by the two
weak-equivalence endpoint projections of its twisted-arrow category.
The unpointed space action
(\cite[Corollary~3.26]{A2CisinskiKan}) identifies constant homotopy
colimits with the classifying-space action, so these projections preserve
the constant-Hom tests by the Kan adjunction. Their natural squares for
$f$ then compare $f$ and $f^{\op}$.
Thus $W_{\mathbb D^{\op}}=W_{\mathbb D}$, completing the dual case.
\end{proof}

\begin{claim}{A2.10}{Proved}{Restricted coefficients have the full strength in the pointed case}
\label{A2-pointed-coefficients}
For a strong pointed derivator, the constant, locally constant, and arbitrary
coefficient versions of each provisional axiom are equivalent. More precisely,
constant Der~$6'$ already implies that every cocartesian square is cartesian;
constant Der~6 already implies the reverse implication. These statements hold
for bare strong derivators, without an enhancement assumption.
\end{claim}
\begin{proof}
We prove the assertion for Der~$6'$; passage to the opposite derivator proves
the other assertion. By A2.1 the constant hypothesis holds in every shift.

First consider the span of finite discrete categories
\[
 \{a_1,a_2\}\ \longleftarrow\
 \{e_1,e_2,e_3\}\ \longrightarrow\ \{b_1,b_2\},
 \qquad
 (e_1,e_2,e_3)\longmapsto(a_1,a_2,a_2),\ (b_1,b_1,b_2).
 \tag{A2.10.1}
\]
Its Grothendieck construction has classifying space the path
$a_1-b_1-a_2-b_2$, with the three edges subdivided; it is contractible.
Use its split (op)fibration, in the variance specified at the start of this
subsection, in the constant-coefficient test. For every $M$ the resulting
square is therefore
\[
 \begin{tikzcd}
 M\amalg M\amalg M
   \ar[r,"h"]\ar[d,"v"']&
 M\amalg M\ar[d,"\nabla"]\\
 M\amalg M\ar[r,"\nabla"']&M ,
 \end{tikzcd}
 \qquad
 h=[i_1,i_2,i_2],\quad v=[i_1,i_1,i_2].
 \tag{A2.10.2}
\]
It is cocartesian by the constant-hocolimit formula, and cartesian by the
hypothesis. The path contraction identifies its last vertex with $M$ and
its structure maps with the folds; these identifications are coherent and
natural in $M$.

There is a specified coherent map from the constant cospan
$(M\xrightarrow{\id}M\xleftarrow{\id}M)$ to
$(M\amalg M\xrightarrow{\nabla}M\xleftarrow{\nabla}M\amalg M)$:
choose $i_1$ at the first endpoint and $i_2$ at the second, and the identity
at the middle. This is induced by a natural transformation of diagrams of
finite sets; both composites to the middle are literally the identity.
Applying the right Kan extension which forms the pullback, and using
(A2.10.2), gives a natural map
\[
 d_M:M\longrightarrow M\amalg M\amalg M,\qquad
 h d_M=i_1,\quad v d_M=i_2.
 \tag{A2.10.3}
\]
Thus its construction does not appeal to an ordinary universal property of
a homotopy pullback in $\mathbb D(e)$.

Since the derivator is pointed, collapse the middle summand in (A2.10.3)
to zero. The resulting natural map
$\delta_M:M\to M\amalg M$ satisfies
$p_1\delta_M=p_2\delta_M=\id$, where $p_1=[\id,0]$ and
$p_2=[0,\id]$. A natural comultiplication with these two counits forces
the canonical map $c:M\amalg N\to M\times N$ to be invertible.
Indeed its inverse is
\[
 \sigma=(\operatorname{pr}_M\amalg\operatorname{pr}_N)
                   \delta_{M\times N}.
 \tag{A2.10.4}
\]
The counits give $c\sigma=\id$; naturality of $\delta$ at the two canonical
maps $M,N\to M\times N$ gives $\sigma c=\id$ on both coproduct summands.
Hence finite products and coproducts are biproducts. With the resulting
commutative-monoid addition on Hom sets, write
$d_M=(\id,\nu_M,\id)$. Equations (A2.10.3) give
$\id+\nu_M=0=\nu_M+\id$, so all Hom monoids are groups.
The same construction in every shift proves that every $\mathbb D(I)$ is
additive.

Next use the span of discrete categories $e\leftarrow\{0,1\}\to e$.
The constant test says that the cocartesian square with upper-left entry
$M\oplus M$ and both outgoing maps $[1,1]:M\oplus M\to M$ is cartesian.
The change of basis
\[
 T=\begin{pmatrix}1&-1\\0&1\end{pmatrix},
 \qquad [1,1]T=[1,0],
 \tag{A2.10.5}
\]
identifies this span with the direct sum of the constant $M$-span and
$0\leftarrow M\to0$. Here is a precise coherence check: first apply
strongness only to the two coherent arrows $[1,1]$ and $[1,0]$, lifting
the isomorphism with components $T$ and $\id$ to $\mathbb D([1])$.
Der~2 makes this lift an isomorphism. Pull it back along the functor
from the span to $[1]$ which identifies its two outgoing edges.
The projection arrow is itself the direct sum of $\id_M$ and $M\to0$.
Left Kan extension preserves this direct sum. Thus the suspension square
is a retract of the tested cartesian square. Retracts of cartesian
squares are cartesian, by their corner-Kan comparison. It follows that
$M\to\Omega\Sigma M$ is invertible, in every shift.

For completeness, the passage from these suspension squares to all
cocartesian squares is a derivator argument. The canonical adjunction
$\operatorname{cof}\dashv\operatorname{fib}$ on $\mathbb D([1])$ and the
natural identity
$\operatorname{cof}^{\,3}\cong\Sigma_{\mathbb D_{[1]}}$ are
\cite[Lemmas~5.12--5.13]{A2GrothPontoShulman}; neither assumes stability.
The suspension is fully faithful because its adjunction unit is invertible.
If an endofunctor $G$ has fully faithful $G^3$, then $G$ is fully faithful:
faithfulness follows immediately, and one lifts $G^2u$ through $G^3$
and cancels the faithful $G^2$ to prove fullness. Consequently
$\operatorname{cof}$ is fully faithful. Its unit
$\id\to\operatorname{fib}\operatorname{cof}$ is invertible, which says
precisely that every cofiber square is cartesian.

Finally append a cofiber square beneath an arbitrary cocartesian square.
Both the lower square and the outer rectangle are cofiber squares, by
the cocartesian pasting lemma. They are therefore cartesian; the
cartesian pasting lemma implies that the original upper square is
cartesian. This proves the unrestricted half-axiom. The implications
from arbitrary to locally constant to constant coefficients then give all
the stated equivalences. The finite incidence and matrix data in
(A2.10.1)--(A2.10.5) are checked integrally by the A2 verification program.
\end{proof}

\begin{claim}{A2.11}{Proved}{The constant axiom makes the pointed envelope one-sided stable}
\label{A2-constant-envelope}
Let $\mathbb D$ be a strong derivator, without a pointedness assumption.
If it satisfies constant Der~$6'$, then every cocartesian square of
$\mathbb D_*$ is cartesian. Dually, constant Der~6 makes every cartesian
square of the copointed envelope $(\mathbb D^{\op})_*^{\op}$ cocartesian.
This conclusion does not yet assert the corresponding implication in
$\mathbb D$ itself.
\end{claim}
\begin{proof}
Put $\mathbb E=\mathbb D_*$, with free pointing $F$ and forgetful $U$.
All the following constructions commute with shifts. The constant tree
test still gives (A2.10.3) in the unpointed derivator. For a pointed
object $N=(s:1\to M)$, compose $d_M$ with
\[
 M\amalg M\amalg M\longrightarrow U(N\vee N\vee N),
 \tag{A2.11.1}
\]
where $\vee$ denotes the coproduct in $\mathbb E$. This is a pointed
map: naturality of $d$ at the coherent section $s$ and
$\nabla_3d_1=\id_1$ identify its value on that section with the common
basepoint. The two fold equations (A2.10.3) descend to pointed wedges.
The counital-comultiplication and middle-component arguments in A2.10
therefore prove that every $\mathbb E(I)$ is additive. No inheritance of
constant left-Kan tests from $\mathbb D$ to arbitrary pointed coefficients
has been assumed.

Apply the unpointed constant free-circle test to $M$. Choose the first
edge of its indexing span and the section $s$ to point that span.
Its pointing extends coherently to the completed square: the colimit of
the constant source-$1$ diagram over a span is $1$. Since $U$ creates
pushouts and pullbacks, this lifted square in $\mathbb E$ is bicartesian.
Its upper-left vertex, pointed on the first copy of $M$, is canonically
\[
 N\vee FM=N\oplus FM;
 \tag{A2.11.2}
\]
indeed its underlying pushout is
$M\amalg_1(1\amalg M)\cong M\amalg M$.
Both outgoing arrows are $[\id,\epsilon_N]:N\oplus FM\to N$.
The invertible triangular matrix
$\left(\begin{smallmatrix}\id&-\epsilon_N\\0&\id\end{smallmatrix}\right)$
turns them into the projection onto $N$. Strongness applied to this one
arrow and then restriction along the span-to-arrow functor gives precisely
the coherence used in A2.10. Consequently the suspension square of
$FM$ is a retract of the lifted cartesian square. Thus
$FM\to\Omega\Sigma FM$ is invertible.

Every $N$ is a retract of $FM$ in the present additive envelope.
The cofiber square (A2.4.4) exists without stability:
$F(s):F1\to FM$ has cofiber $\epsilon_N:FM\to N$ and retraction
$F(M\to1)$. We verify the splitting without invoking a triangulated
structure. In any additive strong pointed derivator, let $i:A\to B$ have
retraction $r$, and let $p:B\to C$ be its coherent cofiber. Cofiber Hom
exactness, proved by the strongness argument in (A2.9.2), factors
$\id_B-ir$ as $qp$ for some $q:C\to B$. The preceding part of the
Puppe Hom sequence is
\[
 \Hom(\Sigma B,T)\longrightarrow\Hom(\Sigma A,T)
 \longrightarrow\Hom(C,T)\xrightarrow{p^*}\Hom(B,T).
 \tag{A2.11.3}
\]
Its first arrow is surjective, since $\Sigma i$ has retraction $\Sigma r$.
Thus $p^*$ is injective. From $pqp=p$ and $rqp=0$ we obtain
$pq=\id_C$ and $rq=0$. The same injectivity makes $q$ unique, hence
natural for coherent split cofiber sequences. This proves the desired
retract, in every shift.

The suspension unit is natural; its value on a retract of $FM$ is a
retract of its invertible value on $FM$. Hence $\id\to\Omega\Sigma$ is
invertible throughout $\mathbb E$. The $\operatorname{cof}^{\,3}$
argument and the pasting argument at the end of A2.10 now prove the
claimed one-sided square property for $\mathbb E$.
\end{proof}

\paragraph{The affine comparison and its obstruction.}
A2.11 alone does not transfer the half-axiom back to the unpointed
derivator; the counterexample A2.12 below shows that such a transfer is
false for constant coefficients. In the enhanced setting,
the canonical functor
\[
 K:\mathcal C\longrightarrow(\mathcal C_*)_{/F1},
 \qquad X\longmapsto(FX\xrightarrow{F(!)}F1)
 \tag{A2.11.4}
\]
has right adjoint $H(q:Z\to F1)=UZ\times_{UF1}1$, with the second map
$1\to UF1=1\amalg1$ selecting the second summand. Indeed the mapping-space
fiber defining maps over $F1$, followed by $F\dashv U$, identifies
$\Map(KX,q)$ with $\Map(X,Hq)$. Thus $K$ is already cocontinuous;
its full faithfulness would suffice to transfer the half-axiom, and
essential surjectivity is unnecessary. Its unit is the comparison
\[
 X\longrightarrow (1\amalg X)\times_{1\amalg1}1.
 \tag{A2.11.5}
\]
The square defining (A2.11.5) is a nonconstant pushout, so assuming it
cartesian would assume an additional implication. In A2.12 its comparison
at $M_0$ is exactly the noninvertible label-changing map $M_0\to M_1$.
For a bare derivator one must in addition avoid the invalid
ordinary-slice substitution discussed below. A2.13 supplies a different
transfer, valid for locally constant coefficients and requiring no slice.

\begin{claim}{A2.12}{Refuted}{Constant coefficients need not give locally constant coefficients}
\label{A2-constant-counterexample}
There is a strong derivator on all small categories satisfying constant
Der~$6'$ and failing its locally constant version, hence also its arbitrary
version. It comes from an explicitly described complete and cocomplete
accessible infinity category. The failing locally constant test has five
objects and seven arrows, including identities. Passing to the opposite
derivator gives the same strict separation for Der~6.
\end{claim}
\begin{proof}
Let
\[
 \mathcal A=D(\Z)_{\ge0}/\Z[0],\qquad P=\{0<1\}.
 \tag{A2.12.1}
\]
Write $a_0=(0\to\Z[0])$ and $a_1=\id_{\Z[0]}$ for its initial and terminal
objects. For $q:V\to\Z[0]$ put
$z(q)=0$ if $H_0(q)=0$, and $z(q)=1$ otherwise.
The mapping space from a connective complex to $\Z[0]$ is discrete,
with set of components $\Hom_{\Ab}(H_0V,\Z)$, so this is an invariant of
the object $q$, not of a chosen chain representative. A morphism $q\to q'$
with $z(q')=0$ forces $z(q)=0$.

Define $\mathcal C$ as the full infinity subcategory of $\mathcal A\times P$
on pairs
\[
 (q,s)\quad\text{with}\quad s\le z(q).
 \tag{A2.12.2}
\]
Thus every augmentation has a label-$0$ version; a label-$1$ version
is allowed exactly for a nonzero augmentation. Its mapping spaces are
those of $\mathcal A$ when the source label is at most the target label,
and are empty otherwise. In particular the two objects
\[
 M_0=(a_1,0),\qquad M_1=(a_1,1)
 \tag{A2.12.3}
\]
are distinct, and the map $b:M_0\to M_1$ with underlying identity is not
an equivalence.

We give its limits, colimits, and accessibility explicitly. The inclusion
$i:\mathcal C\to\mathcal A\times P$ has right adjoint
\[
 R(q,s)=(q,\min\{s,z(q)\}).
 \tag{A2.12.4}
\]
Indeed a map from an allowed pair to an augmentation with $z(q)=0$
already has source label zero, so lowering the target label makes no
change to its mapping space. Consequently limits in $\mathcal C$ are
$R$ applied to the ambient limits. Colimits are created by $i$:
if an ambient colimit augmentation is zero, every augmentation in its
diagram is zero, since each maps to that colimit. Every label in that
case is zero, and their supremum is zero. The empty colimit also has
label zero. Thus every ambient colimit of an allowed diagram is allowed.
This proves all small limits and colimits.

Here is also a concrete accessibility check. Replace an augmented
connective complex by a nonnegative complex of free abelian groups, and
write it as the filtered union of its bounded finite free cell
subcomplexes, with the induced augmentations. Such subcomplexes are
obtained by closing finitely many chosen free generators under their
differentials; the closure is finite because degrees decrease and every
boundary is a finite sum. These filtered colimits compute the derived
colimit. For a label-$0$ object use label zero on every approximant.
For a label-$1$ object choose one degree-zero element whose augmentation
is a fixed nonzero integer. Subcomplexes containing this witness form
a cofinal system and have nonzero augmentation, so may all be labelled
one. Bounded finite free augmented complexes are compact in $\mathcal A$:
their slice mapping spaces are fibers of the usual compact mapping
spaces over the discrete space of augmentations, and filtered colimits
of spaces commute with finite limits. They remain compact with label
zero in $\mathcal C$. With label one, a filtered target diagram whose
colimit has label one has a cofinal tail of label-one objects, giving
the same conclusion; if its colimit has label zero both mapping sides
are empty. There is a set of these labelled finite free complexes.
They generate every object under filtered colimits, proving accessibility.

Now define the entire derivator, not merely its value at a point:
\[
 \mathbb D_{\mathcal C}(I)=
       \Ho\operatorname{Fun}(NI,\mathcal C)
       \qquad(I\in\Cat).
 \tag{A2.12.5}
\]
Restrictions are precomposition; both Kan extensions are the infinity
categorical Kan extensions. Der~1 follows from the decomposition of a
functor category over a disjoint union, and Der~2 from pointwise detection
of equivalences. The limits and colimits just proved give Der~3, and their
comma-category formulas give Der~4. For Der~5 apply the coherent-arrow
construction in $\operatorname{Fun}(NI,\mathcal C)$: any homotopy-category
arrow has a representative, and a commuting square in that homotopy
category has a chosen homotopy between its composites, which supplies a
map of coherent arrows. This gives essential surjectivity and fullness.
These are exactly the infinity-category derivator constructions in
\cite[Example~2.5 and the discussion of Der~5 on p.~5]{A2GrothPontoShulman}.
Hence (A2.12.5) is a strong derivator on all small categories.

Every pushout square in $\mathcal A$ is a pullback: connective derived
abelian groups have this property as in A2.2, and the slice forgetful
functor creates colimits and weakly contractible-index limits. We verify
the constant provisional axiom for every indexing category. Fix a constant
coefficient $M=(q,s)$. Since $i$ creates colimits, any left Kan value of
this constant coefficient has label $s$ if its comma indexing category is
nonempty, and label zero if it is empty. If $s=0$ all labels are zero.
The tested square is a pushout in $\mathcal A$, hence a pullback there;
applying (A2.12.4) to that pullback still gives label zero, as required.

Suppose $s=1$. For a nonempty indexing category its constant homotopy
colimit has nonzero augmentation: choosing any object gives a map from
$q$ to that colimit, and a nonzero augmentation cannot map to a zero
one. If the upper-left comma category of the tested square is nonempty,
the functors to both adjacent comma categories make those nonempty too.
Their labels are both one, and the ambient pullback has the nonzero
upper-left augmentation; (A2.12.4) therefore keeps label one. If the
upper-left comma category is empty, that augmentation is $a_0$.
The ambient pullback is $a_0$, so (A2.12.4) resets its label to zero,
regardless of the adjacent labels. This proves the comparison is an
equivalence in every case. The proof applies to arbitrary small
(op)fibrations in the provisional axiom, not just finite discrete fibers.
By A2.1 it proves the constant version in every shift as well.

For failure with locally constant coefficients use the split opfibration
over the span whose fibers and transition maps are
\[
 \{a_0',a_1'\}\ \longleftarrow\ \{r\}\
 \longrightarrow\ \{b_0',b_1'\},
 \qquad r\longmapsto a_0',\ b_0'.
 \tag{A2.12.6}
\]
Its total category has the two nonidentity arrows
$r\to a_0'$ and $r\to b_0'$, and two isolated objects $a_1',b_1'$.
Give its three-object connected component the constant value $M_0$
and each isolated object the value $M_1$. Every coefficient arrow is an
identity, so the coefficient is locally constant. The resulting square is
\[
 \begin{tikzcd}
 M_0\ar[r]\ar[d]&M_0\amalg M_1\ar[d]\\
 M_0\amalg M_1\ar[r]&M_0\amalg M_1\amalg M_1.
 \end{tikzcd}
 \tag{A2.12.7}
\]
Its three noninitial corners have label one, whereas its initial corner
has label zero. Its ambient pullback in $\mathcal A$ is $a_1$, with
nonzero augmentation. Thus its pullback in $\mathcal C$ is $M_1$, and the
comparison is $b:M_0\to M_1$, not an equivalence.

The underlying finite calculation is integral. The two maps
$\Z^2\to\Z^3$ include coordinates $(1,2)$ and $(1,3)$; their difference
has matrix
\[
 \begin{pmatrix}1&0&-1&0\\0&1&0&0\\0&0&0&-1\end{pmatrix}.
 \tag{A2.12.8}
\]
It is surjective, as witnessed by a $3$-minor of determinant $-1$,
and its kernel is generated by $(1,0,1,0)$. The induced augmentation
of that kernel is the identity of $\Z$. The A2 checker verifies the
category, its opfibration lifts, all coefficient arrows, this integral
kernel certificate, and the differing labels.
An even smaller failure of the arbitrary version is the pushout
$M_1\leftarrow M_0\to M_1$, whose fourth vertex and pullback are $M_1$.

This also exhibits failure of conservative free pointing:
$1_{\mathcal C}=M_1$ and $1_{\mathcal C}\amalg b$ is an equivalence,
although $b$ is not. It refutes the affine-unit transfer (A2.11.5).
The derivator nevertheless has $W_{\mathbb D_{\mathcal C}}=W_{\Z}$.
For clarity use the adjunction characterization already used in A2.4:
a functor $f:I\to J$ belongs to $W_{\mathbb D}$ precisely when every
canonical constant-homotopy-colimit comparison
\[
 \operatorname{hocolim}_{I}M\longrightarrow
 \operatorname{hocolim}_{J}M
 \tag{A2.12.9}
\]
is invertible, by adjunction and Yoneda in the target coefficient.
The underlying complex of the colimit of $M=(V\to\Z[0],s)$ is
$C_*(BI;\Z)\otimes^{\mathbf L}V$. Taking $M=M_0$ therefore detects
integral homology. Conversely an integral-homology equivalence remains
a quasi-isomorphism after tensoring with every connective $V$, hence
makes the underlying comparison in $\mathcal A$ invertible. Its labels
also agree: a constant colimit has label $s$ for a nonempty index and
zero for an empty index, and integral $H_0$ detects emptiness. Thus
(A2.12.9) is an equivalence in $\mathcal C$ for every $M$, proving the
claimed equality. This counterexample separates coefficient strengths;
it does not provide a non-$RT$-fixed localizer.
\end{proof}

\begin{claim}{A2.13}{Proved}{Locally constant coefficients have the unrestricted strength}
\label{A2-locally-constant-full}
For every strong derivator, including an unpointed one, the locally
constant and arbitrary coefficient versions of each provisional axiom
are equivalent. Consequently the exact comparison is
\[
 \text{arbitrary}\ \Longleftrightarrow\ \text{locally constant}
 \ \Longrightarrow\ \text{constant}.
 \tag{A2.13.1}
\]
The last implication is strict in general by A2.12, and is an equivalence
for pointed strong derivators by A2.10. Either locally constant provisional
axiom already makes $W_{\mathbb D}$ $RT$-fixed, by A2.9.
\end{claim}
\begin{proof}
It suffices to prove the assertion for Der~$6'$. The locally constant
hypothesis includes the constant one, so A2.11 applies to
$\mathbb E=\mathbb D_*$. Let $T=UF=1\amalg(-)$.
Use the five-object opfibration (A2.12.6), with coefficient $X$ on its
three-object component and coefficient $1$ on each isolated vertex.
These coefficients are locally constant. The hypothesis makes the natural
split-coproduct square
\[
 \begin{tikzcd}
 X\ar[r]\ar[d]&X\amalg1\ar[d]\\
 X\amalg1\ar[r]&X\amalg1\amalg1
 \end{tikzcd}
 \tag{A2.13.2}
\]
cartesian, with the two distinct new summands on its right and lower edges.
The construction is coherent and natural in $X$, and is valid in every
shift by A2.1. Thus it reconstructs $X$ as the homotopy pullback of
the indicated $TX\to T^2X\leftarrow TX$.

Let $S$ now be any coherent cocartesian square of $\mathbb D$.
The square $FS$ is cocartesian in $\mathbb E$, hence cartesian by A2.11.
Since $U$ creates these squares, $TS$ is cartesian. Moreover $T$ preserves
pushouts, since $F$ preserves colimits and $U$ creates pushouts.
Apply the same argument to the pushout square $TS$ to conclude that
$T^2S$ is cartesian.
Apply (A2.13.2) to the coefficient $S\in\mathbb D(\square)$.
It expresses $S$ as a pointwise homotopy pullback of the three cartesian
squares $TS,T^2S,TS$. Cartesian squares are closed under such limits:
they are the essential image of the right corner Kan extension, and
right Kan extensions in independent directions commute, by their
adjunctions. Hence $S$ is cartesian. This proves the unrestricted half,
and passage to the opposite derivator proves the other half.
\end{proof}

\begin{claim}{A2.14}{Proved}{Constant coefficients already force $RT$-fixedness in the enhanced case}
\label{A2-constant-enhanced-RT}
Let $\mathcal C$ be a complete and cocomplete infinity category and
$\mathbb D_{\mathcal C}$ its homotopy derivator. If either provisional
axiom holds with constant coefficients, then $W_{\mathbb D_{\mathcal C}}$
is $RT$-fixed. For the pushout-to-pullback half an explicit description is
\[
 W_{\mathbb D_{\mathcal C}}
   =\bigcap_{M\in\mathcal C}W_{\mathbb E_M}
   =R\left(\bigcap_{M\in\mathcal C}L_M\right),
 \quad
 \mathbb E_M=\mathbb D_{(\mathcal C_{/M})_*},
 \quad W_{\mathbb E_M}=R(L_M).
 \tag{A2.14.1}
\]
Here $(\mathcal C_{/M})_*$ is the genuine infinity category of retractive
objects $M\to X\to M$, with a specified composite homotopic to $\id_M$.
No accessibility assumption is needed. This assertion is enhanced:
it does not postulate slices of an arbitrary bare strong derivator.
\end{claim}
\begin{proof}
Assume the constant pushout-to-pullback condition. For every $M$ the
slice $\mathcal C_{/M}$ is complete and cocomplete. Its forgetful functor
creates colimits and limits with weakly contractible indexing nerve,
in particular pullbacks. Thus it inherits the constant condition.
By A2.11 its pointed-envelope derivator $\mathbb E_M$ satisfies the
unrestricted pushout-to-pullback implication. A2.9 consequently supplies
a localizing subcategory $L_M\subseteq\Ho(\Sp)$ with
$W_{\mathbb E_M}=R(L_M)$.

We prove the first equality of (A2.14.1), including empty categories.
Write $U_M:(\mathcal C_{/M})_*\to\mathcal C_{/M}$ and $F_M$ for
forgetful and free pointing. For a retractive object $N$ with underlying
object $U_MN=(X\to M)$, its constant homotopy colimit indexed by $I$
has underlying object in $\mathcal C$
\[
 M\amalg_{\operatorname{hocolim}_I M}
                 \operatorname{hocolim}_I X.
 \tag{A2.14.2}
\]
The pushout supplies the section from $M$, and the augmentation to $M$
is induced by that of $X$. Formula (A2.14.2) follows directly from the universal property of the
colimit under the terminal object of the slice; it is the relative
version of (A2.4.1a). If $f:I\to J$ lies in
$W_{\mathbb D_{\mathcal C}}$, its comparisons on the two constant
coefficients $M$ and $X$ are equivalences. Their pushout comparison in
(A2.14.2) is therefore an equivalence for every $N$. The
constant-hocolimit/Yoneda criterion proves
$f\in W_{\mathbb E_M}$ for every $M$.

Conversely suppose that $f\in W_{\mathbb E_M}$ for every $M$ and first
assume $I$ nonempty. Choose $i\in I$. For each $M$, the canonical
comparison
\[
 \tau_f(M):\operatorname{hocolim}_I M\longrightarrow
                   \operatorname{hocolim}_J M
 \tag{A2.14.3}
\]
is a map of retractive objects over $M$: use the augmentations and the
sections selected by $i$ and $f(i)$. Denote this map in
$(\mathcal C_{/M})_*$ by $n_f(M)$. Cocontinuity of $F_M$ identifies
$F_MU_M n_f(M)$ with the constant-hocolimit comparison in $\mathbb E_M$
for the coefficient $F_M(\id_M)$. It is an equivalence by hypothesis.
The natural splitting established in A2.11 makes every retractive object
$N$ a retract of $F_MU_MN$, naturally in $N$. Hence $n_f(M)$ is a retract
of the equivalence $F_MU_M n_f(M)$ in the ordinary arrow category
of $\mathbb E_M(e)$.
It too is an equivalence. Forgetting shows that (A2.14.3) is invertible
for every $M$, and therefore $f\in W_{\mathbb D_{\mathcal C}}$.
This uses the natural splitting inside each pointed envelope; it does
not assert conservativity of free pointing on the unpointed slice.

If $I=J=\varnothing$ there is nothing to prove. The only remaining case
is $I=\varnothing$, $J\ne\varnothing$. For every $N\in\mathbb E_M(e)$,
the hypothesis identifies $\operatorname{hocolim}_J N$ with the zero
object of $\mathbb E_M$. A chosen vertex of $J$ and the augmentation
make $N$ a retract of this colimit. Thus every $\mathbb E_M$ is trivial.
In particular $F_M(\id_M)$ is zero, so the fold
$M\amalg M\to M$ is an equivalence. Now apply the constant condition
with coefficient $M$ to the span having empty root fiber and one object
on each outer fiber. Its square is
\[
 \begin{tikzcd}
 0_{\mathcal C}\ar[r]\ar[d]&M\ar[d]\\
 M\ar[r]&M\amalg M.
 \end{tikzcd}
 \tag{A2.14.4}
\]
Both maps into the lower-right corner are equivalences, so cartesianness
forces $0_{\mathcal C}\simeq M$. This holds for every $M$, making
$\mathcal C$ trivial. In this case $W_{\mathbb D_{\mathcal C}}$ contains
all functors, as required.

Finally intersections of localizing subcategories are localizing, and
$R$ preserves these intersections: both statements follow directly from
the membership test for the suspension cofiber of a functor. Hence
(A2.14.1) proves $RT$-fixedness. The intersections are understood in the
same universe convention as the coefficient classes and A2.9; no
assertion that they are generated by a set is used. The other provisional
half follows by applying the argument to $\mathcal C^{\op}$ and using
the opposite-invariance comparison proved in A2.9.
\end{proof}

\paragraph{Boundary for a bare-derivator slice statement.}
A2.6 is not asserted for arbitrary strong derivators here. Such an assertion
needs a genuine derived slice of a stable derivator over a nonterminal
object. The ordinary slice $\mathbb E(I)/p_I^*V$ and the ordinary fiber of
$\mathbb E(I\times[1])\to\mathbb E(I)$ do not supply it. For example take
$V=\mathbb S\oplus\Sigma\mathbb S$ and its first-summand inclusion
$q:\mathbb S\to V$. In the genuine spectral slice
\[
 \Map_{\Sp_{/V}}(q,q)
 \simeq\Omega\Map_{\Sp}(\mathbb S,\Sigma\mathbb S),
 \qquad\pi_0\Map_{\Sp_{/V}}(q,q)\cong\Z.
\]
An ordinary fiber of homotopy categories identifies these classes under
loops of $\Map(V,V)$ which move the first summand into the second. It
therefore loses precisely the homotopies that a derived slice must retain.
The complete bare-derivator results A2.4 and A2.7 avoid this problem: the pointing
is over a terminal object and its spectral action is supplied directly by
the derivator theorem cited there.

\paragraph{Scope of the coefficient restrictions.}
The arguments above prove shift invariance separately for each of the
arbitrary, constant, and locally constant versions. A2.10 proves that all
three coefficient versions are equivalent for pointed strong derivators.
For a general unpointed strong derivator, A2.13 identifies the locally constant
and unrestricted versions, while A2.12 strictly separates the constant version.
In particular one must not use the identity
fibration with an arbitrary corner coefficient while claiming to prove a
result from the constant-coefficient version. The coefficient-strength
comparison and the shift question are therefore both settled, with their
different pointed and unpointed scopes stated explicitly. A2.14 additionally
settles constant-only $RT$-fixedness for infinity-category derivators. Its
extension to all bare strong derivators would need genuine relative
pointing over arbitrary coefficients, or a relative Brown construction;
neither is assumed here.
